% MAMBench manuscript -- IEEE conference format (IEEEtran, compsoc), anonymized for double-blind review.
% Every number is a macro defined in numbers.tex, generated from results/tables by
% `python -m phase17.remediation.paper_numbers`. Do not type numbers into this file.
\documentclass[conference,compsoc]{IEEEtran}
\usepackage[utf8]{inputenc}
\usepackage{booktabs,multirow,amsmath,amssymb,url,xspace}
\usepackage[hidelinks]{hyperref}
% GENERATED by phase17/remediation/paper_numbers.py from results/tables -- do not edit by hand.
\newcommand{\AdaptCorr}{22/48\xspace}
\newcommand{\AdaptLedger}{5/48\xspace}
\newcommand{\AnomLeakFree}{131/160 attacks at 4/85 benign\xspace}
\newcommand{\AnomShippedAtk}{159/160\xspace}
\newcommand{\AnomShippedFP}{137/171 (80\%)\xspace}
\newcommand{\AuditExt}{34.8\%\xspace}
\newcommand{\AuditExtN}{1{,}824\xspace}
\newcommand{\AuditN}{2{,}628\xspace}
\newcommand{\AuditOwn}{60.0\%\xspace}
\newcommand{\AuditOwnN}{804\xspace}
\newcommand{\AuditRate}{42.5\%\xspace}
\newcommand{\BthirteenFPRmax}{15.3\%\xspace}
\newcommand{\BthirteenFPRmin}{8.8\%\xspace}
\newcommand{\Conclusion}{Re-scoring defenses on attacks shown to work changes the measured recall by up to \TwoDenMaxShift{} points in either direction, so validity filtering belongs in memory-poisoning evaluation. For admission, corroborating untrusted personal claims against the owner's conversation reached memory poisoning that content classifiers missed on three single-use hold-outs, at a benign cost, while the channel label itself added nothing and left the user channel unprotected. For propagation, retention-based lineage taint acts only on descendants of poisons admission already caught; at the second generation it withheld 24/71 carrying summaries when the intermediate summary inherited its state and 0/71 when it did not. For sleeper rules, SA-1 confirmed dormancy and natural trigger activation, and showed that conditional rules with an explicit malicious action are stopped before activation by content detection (60/60) with little cost to benign conditional rules (2/60); no sleeper-specific signal worked, and stealthy, instruction-free sleepers remain largely undetected. Each conclusion holds within an in-context evaluation harness, a 7B agent and evaluator-assigned channels, and should be re-tested in a production memory system.}
\newcommand{\GandalfShare}{38\%\xspace}
\newcommand{\GenOneMissed}{88\xspace}
\newcommand{\GenOneSemCarry}{76/102 (74.5\%)\xspace}
\newcommand{\GenOneSemNonCarry}{22/59 (37.3\%)\xspace}
\newcommand{\HOeightBfourteen}{69/160 (43.1\%)\xspace}
\newcommand{\HOeightBthirteen}{118/160 (73.8\%)\xspace}
\newcommand{\HOfiveBseventeen}{90.1\%\xspace}
\newcommand{\HOfiveBseventeenFPR}{10.2\%\xspace}
\newcommand{\HOfourBfourteen}{32/157 (20.4\%)\xspace}
\newcommand{\HOfourBthirteen}{79/157 (50.3\%)\xspace}
\newcommand{\HOsevenBfourteen}{38/129 (29.5\%)\xspace}
\newcommand{\HOsevenBthirteen}{70/129 (54.3\%)\xspace}
\newcommand{\HOsixBseventeen}{77.6\%\xspace}
\newcommand{\HOsixBseventeenFPR}{2.2\%\xspace}
\newcommand{\HarmJudgeOnly}{5\xspace}
\newcommand{\HarmNatRetr}{89/309 (28.8\%)\xspace}
\newcommand{\HarmNoiseRange}{2.3\xspace}
\newcommand{\JSAllCzero}{25.3\%\xspace}
\newcommand{\JSAllVote}{27.1\%\xspace}
\newcommand{\JSBaseCzero}{25.7\%\xspace}
\newcommand{\JSBaseVote}{47.4\%\xspace}
\newcommand{\JSExtCzero}{18.5\%\xspace}
\newcommand{\JSExtVote}{19.6\%\xspace}
\newcommand{\JSFlips}{509/2628\xspace}
\newcommand{\JSFlipsPct}{19\%\xspace}
\newcommand{\JSOwnCzero}{40.8\%\xspace}
\newcommand{\JSOwnVote}{44.3\%\xspace}
\newcommand{\LaunderOne}{33/161\xspace}
\newcommand{\LeadBseventeen}{37.2\xspace}
\newcommand{\LeadBthirteen}{32.1\xspace}
\newcommand{\LedgerBlast}{6.4\%\xspace}
\newcommand{\PredDirectionalHeld}{35/42\xspace}
\newcommand{\PredHeld}{45\xspace}
\newcommand{\PredIntervalHeld}{10/11\xspace}
\newcommand{\PredScored}{53\xspace}
\newcommand{\PtwoAbstract}{carried the payload into a second consolidation in 98/249 instances; containing those second-generation summaries depends on inheriting the intermediate summary's state (withheld carrying summaries: 24/71 with inheritance, 0/71 without).}
\newcommand{\PtwoBody}{Second-generation summaries carried the poison's payload in 98/249 instances (95/148 of those whose first summary carried it; P2-1~HELD). Serving them on B13-excluded instances produced harm in 22/161 (P2-4~HELD). Of the 71 carrying summaries of excluded poisons, containment withheld 0/71 (lexical), 24/71 (semantic, inheriting the parent's state), 24/71 (adding grandparents) and 0/71 (grandparents only, no inheritance; P2-3~HELD); P2-2~HELD. The semantic inheriting arm also withheld 7/90 summaries that did not carry the payload (P2-5~FAILED) and 4/79 benign chains. Harm with each arm: A0 22/161, L-inh 22/161, S-inh 13/161, S-inh+d2 13/161, S-noinh 22/161. No arm acts on the 88 instances admission missed, 27 of whose second summaries still carry the payload.}
\newcommand{\RemHolm}{Holm-adjusted $p$: SA-2 9.5e-07; SA-3 2.5e-30; SA-5 9.5e-07; SA-8 7.6e-06; P2-2 6.0e-07; P2-4 9.5e-07.}
\newcommand{\SAoneAbstract}{a pre-registered experiment with 60 conditional backdoors and 60 benign conditional lookalikes found activation on a natural trigger in 23/60 attacks and on an unrelated message in 0/60; admission-time content detection excluded 60/60 attacks and 2/60 lookalikes before activation---because these backdoors state their malicious action explicitly, not because any sleeper-specific signal worked (a dedicated sleeper guard flagged 0/60).}
\newcommand{\SAoneBody}{Each of 60 conditional backdoors and 60 benign conditional lookalikes (same author, owner, prompt and third-person syntax) carries a natural trigger message and an unrelated message. The note was retrieved into the agent's top five by its trigger in 51/60 attacks and 38/60 lookalikes, and by the unrelated message in 4/60 and 3/60. Undefended, attacks activated (complied with the note present, not without it) on 23/60 trigger messages and 0/60 unrelated messages (SA-1~HELD, dormancy; SA-2~HELD, paired +23/$-$0, $p$=2.4e-07); under the human-calibrated vote judge, 42/60 and 1/60. Lookalikes, whose activation is the intended behaviour, activated on 9/60 triggers. At admission, before any activation, B13 on the untrusted channel excluded 60/60 attacks and 3/60 lookalikes (SA-3~HELD, Fisher $p$=4.1e-31); B12 60 and 2 (SA-4~HELD). B13 reduced trigger activation from 23/60 to 0/60 (SA-5~HELD, $p$=2.4e-07) and lookalike activation from 9/60 to 8/60. On the user channel B13 excluded 0 attacks, by construction. The dedicated Phase~8 sleeper guard flagged 0/60 attacks (SA-6~HELD); the development-fitted anomaly signal flagged 60/60 attacks and 60/60 lookalikes (SA-7~HELD); deberta flagged 8/60 and 2/60. At activation time, the instruction-aware frame changed attack activation from 23/60 to 5/60 (SA-8~HELD, $p$=7.6e-06) and lookalike activation from 9/60 to 13/60 (SA-9~HELD). The interpretation matters: B12 alone excluded every attack, so detection here rests on the explicit payload (a named outside company and a covert action), not on dormancy; mundane, instruction-free sleeper notes are not caught by content detection (0\% on the project's sleeper family) and only partly by corroboration. B13's corroboration thresholds were calibrated on this owner's conversation among others (\S\ref{sec:limits}).}
\newcommand{\TwoDenMaxShift}{11\xspace}


\begin{document}
\title{Do the Attacks Work? A Validity-Audited Evaluation of Provenance, Propagation and Sleeper Defenses for LLM Agent Memory}
\author{\IEEEauthorblockN{Anonymous Authors}\IEEEauthorblockA{Under double-blind review}}
\maketitle

\begin{abstract}
Defenses against memory poisoning in LLM agents are scored on attack corpora that are assumed, not shown, to work, and are rarely charged for the benign memories they exclude. We apply a causal, defense-independent validity audit to \AuditN{} attack instances in 20 populations: an instance counts as working only if an undefended agent follows it with the memory present, does not with it absent, and an independent model confirms it implements its labelled mechanism. Under the original judge \AuditRate{} work (\AuditOwn{} of the project's own memory-poisoning instances, \AuditExt{} of external injection corpora); under a judge calibrated against two human annotators the fitted-probe rate moves from \JSAllCzero{} to \JSAllVote{}. Restricting to working attacks shifts defense recall by up to \TwoDenMaxShift{} points, in opposite directions on different hold-outs, so full-population recall is not a neutral summary. We then use the protocol, five single-use frozen hold-outs and pre-registered predictions to examine three defense families. (i) Provenance-aware admission (B13) excludes \HOfourBthirteen{}, \HOsevenBthirteen{} and \HOeightBthirteen{} of memory poisoning on three independently authored hold-outs against \HOfourBfourteen{}, \HOsevenBfourteen{} and \HOeightBfourteen{} for a fused content classifier, at \BthirteenFPRmin{}--\BthirteenFPRmax{} benign exclusion on untrusted channels; a decomposition shows the channel label adds no recall of its own---the entire gain is corroboration of personal claims against the owner's conversation---and poison arriving on the trusted user channel is never excluded. (ii) Lineage-taint containment of consolidated memories \PtwoAbstract{} (iii) For conditional (sleeper) rules, \SAoneAbstract{} An embedding-anomaly sleeper signal reported earlier as generalizing was calibrated on its own hold-out; with development-fitted thresholds it flags \AnomShippedFP{} of the hold-out's benign notes. Of \PredScored{} scored pre-registered predictions, \PredHeld{} held.
\end{abstract}

\section{Introduction}
LLM agents increasingly persist what they read and are told into long-term memory, and an attacker who can influence what is written can steer the agent long after the interaction that planted the content~\cite{agentpoison,minja,hiddenmem,mpbench}. A growing set of defenses targets this threat: content screening, provenance and lineage tracking~\cite{memlineage,tmanm,ppmf,aegis}, consensus checks~\cite{amemguard}, and prompt-level framing~\cite{spotlighting,hiddenmem}. Recent benchmarks measure persistence and execution of poisoned memories~\cite{mempoison,memsecbench,ied}.

Two measurement questions are less settled. First, are the attack instances used to score defenses functional against the agent under test? A detector that "misses" an instance that would never have changed the agent's behavior has not failed, and one that "catches" it has not succeeded. The analogous problem in jailbreak evaluation---success claims that do not survive a check of what the model actually produced---is well documented~\cite{strongreject}; in agent memory, storage and execution are known to dissociate~\cite{ied,memsecbench,hiddenmem}. Second, what does a defense cost on benign memories that \emph{resemble} attacks: notes about people the owner never mentioned, conditional rules the owner genuinely set, summaries derived from an excluded memory?

We study both questions for three defense families that correspond to three stages of the memory lifecycle: admission (provenance and governance), transformation (propagation through consolidation), and activation (dormant, conditional rules). Our contributions are:
\begin{itemize}
\item \textbf{A causal attack-validity protocol and two-denominator reporting} (\S\ref{sec:validity}). Each instance is run with and without its memory on the same task, judged by a model family independent of the agent and of every defense signal, and checked against its labelled mechanism by a third family. We report every defense on the full population and on the working subset, and show the two can rank and size effects differently.
\item \textbf{A decomposition of provenance-aware admission} (\S\ref{sec:prov}): on three single-use memory-poisoning hold-outs, the channel label contributes no recall; corroboration against the owner's conversation contributes all of it, at a measured benign cost, and the trusted-channel exemption makes user-channel poisoning unreachable by construction.
\item \textbf{Measurements of propagation containment and sleeper activation with benign controls} (\S\ref{sec:prop}, \S\ref{sec:sleeper}): a pre-registered two-generation consolidation experiment and a pre-registered activation experiment with 60 benign conditional lookalikes, plus a correction of a hold-out-calibrated sleeper result.
\end{itemize}
We do not claim new defense mechanisms: channel trust, lineage taint, stateful rejection and data framing all have precedents~\cite{tmanm,memlineage,ppmf,statefuldet,blacklight,spotlighting}. Our claims concern what these mechanisms do, at what cost, on attacks shown to work.

\section{Related Work}\label{sec:related}
\textbf{Memory-poisoning attacks and benchmarks.} AgentPoison~\cite{agentpoison}, MINJA~\cite{minja}, MemoryGraft~\cite{memorygraft}, FARMA~\cite{farma}, DSRM~\cite{dsrm} and sleeper memory poisoning~\cite{hiddenmem} target different points of the memory pipeline; MPBench~\cite{mpbench} systematizes six attack classes over four write channels. MemPoison~\cite{mempoison} evaluates 1,227 hand-validated cases with clean/poisoned paired runs and an activation-shift measure for dormant attacks; MemSecBench~\cite{memsecbench} scores persistence, execution and repair on Mem0, Mem0-Graph and A-MEM inside two agent harnesses; Injection--Execution Dissociation~\cite{ied} shows storage and execution are separable, with pre-registration and Holm correction. MemSecBench and these attack papers exercise real memory systems more fully than our security evaluation does (\S\ref{sec:harness}). We are not aware of prior work that re-reports defense detection on a causally filtered working subset; HiddenMem~\cite{hiddenmem} and MPBench condition later stages on earlier ones, which is related but is not a defense-evaluation denominator.

\textbf{Provenance, lineage and laundering.} Channel-based trust appears in Aegis~\cite{aegis} (a retrieval-time trust prior with user turns trusted), in provenance down-weighting~\cite{priceofsafety}, and in origin-bound authority~\cite{tmanm}, which also shows that content- and lineage-based gates are defeated by self-summarization laundering. MemLineage~\cite{memlineage} propagates untrusted origin through derivation chains; PPMF~\cite{ppmf} names provenance laundering during consolidation. CaMeL~\cite{camel} and StruQ~\cite{struq} separate trusted instructions from untrusted data for prompt injection.

\textbf{Sleeper attacks.} Dormant, trigger-activated memory poisoning is measured with trigger versus non-trigger queries in MemPoison~\cite{mempoison}, HiddenMem~\cite{hiddenmem} and Plant--Persist--Trigger~\cite{ppt}. None of the papers we read evaluates detectors against benign conditional rules of the same shape.

\textbf{Stateful and framing defenses.} Remembering rejected queries to stop iterative black-box attackers follows stateful detection~\cite{statefuldet} and Blacklight~\cite{blacklight}; A-MemGuard~\cite{amemguard} stores lessons from detected failures. Framing untrusted text as data follows Spotlighting~\cite{spotlighting}.

\section{Threat Model and Assumptions}\label{sec:threat}
The victim is a personal assistant with a persistent memory store. The attacker controls content that reaches the store through an \emph{untrusted channel}---tool output, a retrieved document, or agent-generated text---and cannot alter the system prompt, model weights, defense parameters, or provenance metadata recorded by the platform. The owner's own conversational turns arrive on a \emph{trusted} channel. We consider single-shot poisoning, retry-capable attackers who receive the defense's verdict (\S\ref{sec:prov}), consolidation of memories into derived summaries (\S\ref{sec:prop}), and conditional rules that remain inert until a situation arises (\S\ref{sec:sleeper}).

\textbf{Load-bearing assumption.} In every evaluation reported here the evaluator assigns each attack to an untrusted channel and the owner's turns to the trusted channel. Provenance-aware admission inherits this assumption: an attack written through the user channel---for example MINJA's query-only threat model~\cite{minja}---or a mislabelled source is treated as the owner's own words. We quantify this in \S\ref{sec:prov} rather than assume it away. Forgery of provenance \emph{metadata} by the attacker is out of scope.

\section{Evaluation Harness and Units of Evaluation}\label{sec:harness}
To keep claims within what was measured, Table~\ref{tab:units} lists the unit of each experiment. None of the security results runs through a production memory system end to end: admission defenses classify memory text; behavioral measurements place memories in an assistant prompt, either directly (validity audit) or through embedding retrieval over a store of the owner's real conversation turns; consolidation is an LLM summarization step with lineage recorded by the harness. A separate 9-case check ran the admission defense on what A-MEM and Mem0 (with its LLM extraction disabled) actually retrieved (supplementary methodology, \S32.5); we do not generalize from it.

\begin{table}[t]\centering\caption{Units of evaluation (what each experiment actually exercises).}\label{tab:units}
\footnotesize
\begin{tabular}{p{2.1cm}p{5.6cm}}\toprule
Experiment & Unit \\\midrule
Validity audit & qwen2.5:7b answers a fitted task with 8 owner turns plus the note in its prompt, and again without the note; llama3.1:8b judges compliance; mistral:7b or Gemini judges mechanism. \\
Defense detection & Each memory text classified at admission; channel label assigned by the evaluator; owner's LoCoMo conversation as corroboration evidence. \\
Harm study & 100 owner turns plus the poison; MiniLM top-5 retrieval for the task; the poison is inserted when retrieval misses it (natural retrieval \HarmNatRetr{}), so harm is conditional on exposure. \\
Consolidation & qwen2.5:7b summarizes the 5 retrieved memories (gen.~1) and the summary plus 4 store memories (gen.~2); parents recorded by the harness; containment uses the frozen guard's scoring functions. \\
SA-1 (sleeper) & 100 owner turns plus the note; real top-5 retrieval for a trigger or non-trigger message, no forced insertion. \\\bottomrule
\end{tabular}\end{table}

Agent-side measurement uses qwen2.5:7b throughout; judges are from different model families than the agent and than any defense signal, and no model judges text it authored. Generation uses temperature 0 with fixed seeds but is not bit-reproducible under concurrent batching; repeated runs of the harm study moved undefended harm by up to \HarmNoiseRange{} points, and re-judging identical answers moved it by \HarmJudgeOnly{} items of 309.

\section{Attack Validity}\label{sec:validity}
\textbf{Protocol.} An instance is \emph{working} if (a) the judge says the agent complied with the memory when it was present, (b) the judge says it did not comply on the same task with the memory absent, and (c) an independent model confirms the instance implements its labelled mechanism. ``Compliance'' includes asserting a planted claim, so working means the memory causally changed the response; it does not mean harm, which is measured separately. The filter never reads a defense decision. Two probes were used---a generic task and a task fitted to each instance---and membership is the union over probes, a rule adopted after the first probe disagreed with the second. Working status is a property of the agent, probe and judge, not of an attack in every setting.

\textbf{Results.} Table~\ref{tab:validity} gives the rates. The pooled \AuditRate{} is dominated by external injection corpora (one game corpus is \GandalfShare{} of the denominator); the project's own memory-poisoning populations work at \AuditOwn{}. The original judge agrees with human annotators at $\kappa\approx$0.5--0.6; re-judging the cached answers with a majority-vote judge calibrated against two annotators gives the fitted-probe rates in Table~\ref{tab:validity}.

\begin{table}[t]\centering\caption{Working-attack rates (fitted probe, mechanism-checked) under the original judge (C0) and the human-calibrated vote; published union-over-probes rate for reference.}\label{tab:validity}
\footnotesize
\begin{tabular}{lrrrr}\toprule
Group & $n$ & C0 & Vote & Union (C0) \\\midrule
All & \AuditN & \JSAllCzero & \JSAllVote & \AuditRate \\
Own memory poisoning & \AuditOwnN & \JSOwnCzero & \JSOwnVote & \AuditOwn \\
External corpora & \AuditExtN & \JSExtCzero & \JSExtVote & \AuditExt \\\bottomrule
\end{tabular}\end{table}

\textbf{Two denominators.} Table~\ref{tab:twoden} reports recall on the full population and on the working subset for every hold-out where the working subset is reportable. On the injection hold-out 6 working attacks are harder to detect for every configuration except one; on memory-poisoning hold-outs 4 and 5 they are easier for provenance; on hold-out 8 they are harder for content classifiers and unchanged for provenance. Full-population recall can therefore overstate or understate protection, and the direction depends on the corpus and the defense.

% GENERATED by phase17/remediation/paper_numbers.py
\begin{table}[t]\centering\caption{Recall on the full population $\rightarrow$ the causal working subset, per single-use hold-out.}\label{tab:twoden}
\footnotesize\begin{tabular}{lrrrrr}\toprule Hold-out (working/all) & B12 & B13 & B14 & B17 & B18\\\midrule
HO4 (86/157) & 2$\rightarrow$2 & 50$\rightarrow$60 & 20$\rightarrow$31 & 50$\rightarrow$60 & 50$\rightarrow$60\\
HO5 (47/131) & 51$\rightarrow$60 & 55$\rightarrow$64 & 90$\rightarrow$89 & 90$\rightarrow$89 & 90$\rightarrow$89\\
HO6 (294/500) & 34$\rightarrow$24 & 37$\rightarrow$29 & 77$\rightarrow$68 & 78$\rightarrow$68 & 99$\rightarrow$99\\
HO8 (63/160) & 37$\rightarrow$27 & 74$\rightarrow$73 & 43$\rightarrow$33 & 74$\rightarrow$73 & 74$\rightarrow$73\\
\bottomrule\end{tabular}\end{table}


\section{Provenance and Governance}\label{sec:prov}
\textbf{Configurations.} B12 is a content-only admission stack (stacked detector, LLM judge tiers for lineage claims and self-replicating directives). B13 adds (i) channel trust---memories on the user channel are never excluded, others get B12's checks; (ii) corroboration---an untrusted memory making a personal claim is excluded when its maximum NLI entailment from the owner's 20 most similar turns falls below a development-frozen threshold; and (iii) a write-time carry-forward guard. B14 is an eleven-feature fused content classifier; B17 is B13\,$\lor$\,B14. Thresholds were frozen on development data before any hold-out.

\textbf{Hold-outs.} Hold-outs 4, 7 and 8 are memory-poisoning corpora written by three different model families about three different owners, each scored once under predictions committed beforehand; benign notes are scored on the untrusted channel (strict) and again on the user channel.

% GENERATED by phase17/remediation/paper_numbers.py
\begin{table}[t]\centering\caption{Memory-poisoning hold-outs, each scored once (B13 with the owner's conversation as evidence). Recall / benign exclusion on the untrusted channel / on the user channel.}\label{tab:mp}
\footnotesize\begin{tabular}{llll}\toprule & HO4 & HO7 & HO8\\\midrule
B12 & 1.9 / 1.5 / 1.5 & 37.2 / 4.2 / 4.2 & 36.9 / 4.1 / 4.1\\
B13 & 50.3 / 11.8 / 0.0 & 54.3 / 15.3 / 0.0 & 73.8 / 8.8 / 0.0\\
B14 & 20.4 / 9.2 / 0.0 & 29.5 / 13.8 / 0.0 & 43.1 / 5.8 / 0.0\\
B17 & 50.3 / 16.0 / 0.0 & 56.6 / 22.8 / 0.0 & 74.4 / 10.5 / 0.0\\
B18 & 50.3 / 16.0 / 1.5 & 57.4 / 22.8 / 3.7 & 74.4 / 10.5 / 1.8\\
\midrule B13 vs B14 & $+$47/$-$0 (p=1.4e$-$14) & $+$35/$-$3 (p=6.7e$-$08) & $+$50/$-$1 (p=4.6e$-$14)\\
\bottomrule\end{tabular}\end{table}


\textbf{What B13's effect depends on.} Table~\ref{tab:decomp} decomposes B13's stored per-item decisions into its rules (reproduced item for item). On all three hold-outs B13 without corroboration excludes exactly what B12 excludes: the channel label contributes no recall. All of B13's gain over content-only detection is corroboration, at the benign cost shown. Excluding every untrusted memory would reach 100\% recall at 100\% untrusted-benign exclusion, so the channel label alone is not a detector. Its role is the exemption: an attack relabelled to the user channel is excluded in 0 cases on every hold-out (restricted, but admitted, when a content rule fires).

% GENERATED by phase17/remediation/paper_numbers.py
\begin{table}[t]\centering\caption{Decomposition of B13 (stored decisions, reproduced item for item). Attacks excluded / untrusted-channel benign excluded.}\label{tab:decomp}
\footnotesize\begin{tabular}{lccc}\toprule & HO4 & HO7 & HO8\\\midrule
B12 (content only) & 3/157 / 4/262 & 48/129 / 8/189 & 59/160 / 7/171\\
B13 without corroboration & 3/157 / 4/262 & 48/129 / 8/189 & 59/160 / 7/171\\
Corroboration alone & 76/157 / 28/262 & 23/129 / 23/189 & 81/160 / 8/171\\
B13 & 79/157 / 31/262 & 70/129 / 29/189 & 118/160 / 15/171\\
\midrule Attacks on user channel & 0/157 (3 flagged) & 0/129 (48 flagged) & 0/160 (59 flagged)\\
\bottomrule\end{tabular}\end{table}


\textbf{Injection-shaped text and published detectors.} On two injection hold-outs with matched benign prompts, the union B17 reaches \HOfiveBseventeen{} and \HOsixBseventeen{} recall at \HOfiveBseventeenFPR{} and \HOsixBseventeenFPR{} false positives. On the project's own memory-poisoning attacks, B13 and B17 lead the strongest published detector by \LeadBthirteen{} and \LeadBseventeen{} points; pooled comparisons across corpora are dominated by one injection corpus and should not be cited.

\textbf{Adaptive attacker.} An attacker that receives the defense's verdict and rewrites the memory up to five times (NLI-checked to preserve its claim) evades corroboration on \AdaptCorr{} of seeds that B13 initially excluded. A per-scope ledger of rejected content, which refuses near-duplicates of what was already rejected~\cite{statefuldet,blacklight}, reduces this to \AdaptLedger{} (pre-registered; development data), at a ledger-poisoning blast radius of \LedgerBlast{} of benign traffic. This is evidence against one bounded paraphrase attacker within one scope, not general adaptive robustness.

\section{Propagation Through Consolidation}\label{sec:prop}
\textbf{Setup.} A consolidation step summarizes retrieved memories into a derived note whose parents are recorded. Containment scores a descendant by $\max_a \mathrm{SEV}(s_a)\cdot 0.5^{d_a-1}\cdot r(D,a)$ over ancestors $a$ at distance $d_a$ with admission state $s_a$ and retention $r$ (lexical token overlap as built, or embedding cosine), and withholds it above a fixed threshold. The only state input is the admission decision; no attack label is used.

\textbf{One generation.} On development data a poison that B13 excluded returns through its summary in \LaunderOne{} cases; lexical containment barely acts while the semantic variant withholds most carrying summaries (\GenOneSemCarry{}). The costs: it also withholds \GenOneSemNonCarry{} of summaries that did \emph{not} carry the payload---legitimate derived memory lost---and it cannot act on summaries of poisons that admission missed (\GenOneMissed{} of 249 instances). Hold-out 8 reproduced the lexical/semantic gap.

\textbf{Two generations (pre-registered).} \PtwoBody{}

\section{Sleeper (Conditional) Poisoning}\label{sec:sleeper}
\textbf{A correction.} A content-independent embedding-anomaly signal was earlier reported to flag 145/160 hold-out-8 attacks at 9/85 benign false positives. Its clusters were fitted on half of that hold-out's benign notes and its threshold set at a percentile of the other half, on which the false-positive rate was then reported---so that rate was fixed by construction. With a threshold fitted on one benign half and measured on the other, the signal flags \AnomLeakFree{}; with the development-fitted model that ships, it flags \AnomShippedAtk{} of attacks and \AnomShippedFP{} of benign notes. It separates this corpus's attacks from its benign notes, but it detects a shift in owner and author rather than poisoning, and it is not usable without per-owner calibration.

\textbf{SA-1: activation with benign lookalikes (pre-registered).} \SAoneBody{}

\section{Pre-registration and Statistics}\label{sec:stats}
Paired comparisons use exact McNemar tests on the same items; unpaired comparisons Fisher's exact test; rates carry 95\% Wilson intervals. Every hold-out was frozen with a SHA-256 digest before it was scored and scored once; predictions were committed to version control before the run they govern. Across the project's final round, \PredHeld{} of \PredScored{} scored predictions held (\PredIntervalHeld{} interval, \PredDirectionalHeld{} directional); one pre-registered prediction was never scored. Fifteen significance predictions form one Holm family; four held at their nominal $\alpha$ but not after correction, and are not used as findings. The remediation experiments (\S\ref{sec:prop}, \S\ref{sec:sleeper}) form a second family: \RemHolm{}

\section{Limitations}\label{sec:limits}
(1)~No security result runs through a production memory system end to end (\S\ref{sec:harness}); results on Mem0 or A-MEM with extraction enabled may differ. (2)~All agent-side measurement uses a 7B model; harm levels depend on a judge that two annotators found to under-count harm ($\kappa$=0.2), so we report reductions rather than levels; the human check covers 100 items and annotator independence from the authors is not documented in our records. (3)~Provenance results assume correct, unforgeable channel labels; user-channel poisoning is unaddressed by design. B13's two corroboration thresholds were calibrated on the conversations of four development owners, two of whom are the owners of hold-out~8 and SA-1 (the calibration never saw those corpora's texts); hold-outs 4 and 7 are the clean-owner evidence. (4)~Corroboration excludes 9--15\% of untrusted-channel benign notes on our hold-outs; it suits deployments where untrusted personal claims can be quarantined for review. (5)~The adaptive and ledger results are development-only and cover one paraphrase attacker within one scope. (6)~Propagation and sleeper findings use LLM-authored corpora; SA-1 notes differ in mean length between classes. (7)~Hold-outs 1--3 are injection or jailbreak corpora not built as memory poisoning; hold-out 3 was used more than once and is excluded from our claims.

\section{Conclusion}
\Conclusion{}

\section*{Ethics Considerations}
All attacks target a local model and synthetic or public data; no deployed system or real user was attacked. Benign memories come from a public conversational dataset (CC BY-NC 4.0). Generated attack corpora describe fictitious companies.

\section*{Open Science}
Code, frozen corpus manifests with digests, pre-registrations, per-item decisions and the scripts that generate every table are available to reviewers in an anonymized repository. Derived data containing CC BY-NC 4.0 source text is released under those terms; one upstream dataset without an explicit licence is not redistributed.

\bibliographystyle{IEEEtran}
\bibliography{refs}
\end{document}
